\documentclass[11pt]{article}
\usepackage[margin=1in]{geometry}
\usepackage{amsmath,amssymb}
\title{A holomorphic lattice VOA of charge eight and weight-one dimension 248}
\author{The Last Math Competition, conjecture 00000004632}
\date{}
\begin{document}
\maketitle

The lattice vertex operator algebra \(V_{E_8}\) is holomorphic, has central
charge \(8\), and has weight-one dimension \(248\). This contradicts the
conjecture's restriction to weight-one dimensions \(0\) and \(1\) below
central charge \(24\).

\paragraph{The lattice.}
Take \(D_8=\{x\in\mathbb Z^8:\sum_i x_i\in2\mathbb Z\}\), let
\(h=\tfrac12(1,1,1,1,1,1,1,1)\), and set \(L=D_8\cup(h+D_8)\).
This is the standard Euclidean coordinate realization of \(E_8\).
The index of \(D_8\) in \(\mathbb Z^8\) is two, and \(2h\in D_8\) while
\(h\notin D_8\). Thus \(L\) is a lattice containing \(D_8\) with index two,
and its covolume is one. For \(x\in D_8\),
\(\|x\|^2\equiv\sum_i x_i\equiv0\pmod{2}\). Also
\(\|x+h\|^2=\|x\|^2+\sum_i x_i+2\) is even. Polarization now gives integral
inner products in \(L\). Hence \(L\) is positive definite, even and
unimodular, of rank eight.

\paragraph{The weight-one space.}
We use the standard lattice-VOA construction for positive-definite even
lattices \cite{dongmason}. For a lattice \(L\) of rank \(r\), it gives a VOA
\(V_L\) of central charge \(r\). Its simple modules are indexed by
\(L^*/L\), so \(V_L\) is holomorphic when \(L\) is unimodular. Its graded
space is the tensor product of the Heisenberg oscillator space and the
twisted group algebra of \(L\), with a basis vector \(e^\alpha\) of weight
\(\|\alpha\|^2/2\) for each \(\alpha\in L\). The weight-one vectors thus come
from the \(r\) first oscillators over the vacuum or one vector \(e^\alpha\)
for each root \(\alpha\) of squared norm two.

There are \(4\binom82=112\) integral roots: they have two entries \(\pm1\)
and six zero entries. For a half-integral root, all eight entries are
\(\pm\tfrac12\), since even one entry of absolute value at least
\(\tfrac32\) would give squared norm greater than two. If \(m\) entries are
negative, their coordinate sum is \(4-m\), which is even precisely when
\(m\) is even. Flipping the first sign pairs every even sign pattern with
an odd one, so there are \(2^7=128\) half-integral roots. Consequently
\[
 c(V_L)=8,\qquad \dim (V_L)_1=8+112+128=248.
\]

\paragraph{Formal verification.}
The Lean project doubles every Euclidean coordinate. Its lattice
membership predicate requires eight coordinates of one parity and sum
divisible by four. It verifies evenness and spans the lattice with the
doubled rows \(4e_0,2(e_1-e_0),\ldots,2(e_6-e_0),(1,\ldots,1)\).
These rows give a lower-triangular matrix of determinant \(256\), which
becomes determinant one on dividing each row by two. The formal root
enumerator lists exactly the lattice vectors of doubled squared norm
eight, and proves that the entries are distinct. It verifies the counts
\(112\) and \(128\) and computes \((c,\dim V_1)=(8,248)\) from the rank
and root count. The lattice-VOA theorem above identifies those data with
the central charge and weight-one dimension of the holomorphic VOA.

\begin{thebibliography}{1}
\bibitem{dongmason}
C.~Dong and G.~Mason, \emph{Holomorphic vertex operator algebras of small
central charge}, Pacific J. Math. \textbf{213} (2004), 253--266.
\end{thebibliography}
\end{document}
